\runninghead{The Clay and the Faithful Potter}
\properlane{clay-and-faithful-potter}{The Clay and the Faithful Potter: Our Failure and God's Fidelity}
\label{sec:clay-and-faithful-potter}

The people who pray Isaiah's cry say of themselves that they have strayed,
that their hearts are hard, that their justice is a soiled rag, that no one calls on God's name. Paul's Corinthians,
whom he thanks God for, are the church he will spend the rest of the letter
rebuking. The Gospel tells servants to be found awake. On what ground can such
people hope to be found blameless on the day? The prophet's people name their
ground in the last verse of the reading: ``And now, O Lord, thou art our father, and we
are clay: and thou art our maker, and we all are the works of thy hands''
(64:7; Douay 64:8). Jerome and Thomas Aquinas, commenting on the passage,
refuse to make God the author of the people's sin and rest the people's hope
of salvation on God's mercy. Augustine, on two of the Mass's psalms, says that
the turning back and the good fruit are God's gifts. John Chrysostom and
Aquinas, on the second reading, find the ground of Paul's confidence in the
sentence ``God is faithful''. Isaiah's people appeal to God as Father and
maker, Paul to God as faithful to his call: the hope of those who wait rests
on who God is, not, as Jerome puts it, on their merits.

\subsection{Jerome and Aquinas: why hast thou made us to err?}

The reading opens with a question that sounds like an accusation. In the Latin
the people ask, \latin{Quare errare nos fecisti, Domine, de viis tuis;
indurasti cor nostrum ne timeremus te?}, ``Why hast thou made us to err, O
Lord, from thy ways: why hast thou hardened our heart, that we should not
fear thee?'' (63:17). Both commentators refuse to read it as a statement that
God causes sin, and they give different reasons.

Jerome says that God is not the cause of error and hardness, but that his
patience seems to be. He waits for the sinner's salvation and does not
correct him at once, and the sinner, left uncorrected, goes further. He gives
the example of Pharaoh, whose heart God is said to harden in Exodus and whom
Paul uses in Romans~9, and the example of a physician who hides the knife from
a patient until the moment of cutting: the delay is for the patient's sake,
though it looks like neglect (\work{Commentarii in Isaiam} XVII, PL~24,
cols.~643--644). Aquinas puts it more briefly. God made them err
\latin{permittendo, non immittens gratiam, et non corrigens}, ``by permitting
it, not sending grace and not correcting''; and \latin{indurasti} means
\latin{permisisti}, ``thou hast hardened'' means ``thou hast permitted''
(\work{Expositio super Isaiam} c.~63). The English proclaimed in the United
States asks why God \emph{lets} his people wander, a permissive rendering
close to Aquinas's \latin{permittendo}. The hard form is the Latin's, and it
is the form the Fathers explained.

The question ends in a petition: \latin{Convertere propter servos tuos},
``return for the sake of thy servants''. In the Latin it asks God himself to
turn back. Jerome reads it two ways: \latin{Convertere, ait, Domine, sive
converte nos propter servos tuos}, ``Turn, he says, O Lord, or, turn us, for
thy servants' sake''. The second reading he develops with a clause the Greek
version adds, ``that we may possess a little of thy holy mountain'': having a
small part, we may come to the whole and \latin{contemplemur gloriam
Unigeniti}, ``behold the glory of the Only-begotten'' (cols.~643--644).
\latin{Converte nos}, ``turn us'', is Jerome's reading of the petition, not the
wording of the verse in any of the texts he translates; it is the wording of
the psalm's refrain. Aquinas brings the two petitions together in another way. He takes
\latin{convertere} as a plea for God to turn \latin{de ira ad misericordiam},
``from anger to mercy'', and sets beside it two other verses: a psalm that
asks God to return (\latin{convertere, Domine}, Ps 90 (89):13), and the end
of Lamentations, \latin{converte nos, Domine, et convertemur}, ``convert us,
O Lord, and we shall be converted'' (Lam 5:21).

\subsection{Jerome and Aquinas: ``an inverted order''}

The second hard saying comes in the confession: \latin{Ecce tu iratus es, et
peccavimus}, ``behold thou art angry, and we have sinned'' (64:4; Douay
64:5). Read as it stands, the anger comes first and the sin follows it. Jerome
and Aquinas give the same answer in nearly the same words. Jerome: \latin{Ordo
praeposterus. Non quia tu iratus es, nos peccavimus; sed quia nos peccavimus,
tu iratus es}, ``The order is inverted. It is not because thou wast angry that
we sinned, but because we sinned that thou wast angry'' (col.~647). Aquinas:
\latin{praeposterus ordo: ideo iratus, quia peccavimus: et tamen salvabimur,
tua misericordia}, ``an inverted order: thou wast angry because we sinned;
and yet we shall be saved, by thy mercy'' (c.~64). Both then read the end of
the verse, ``and we shall be saved'', as the people's hope, and both place the
hope in mercy. Jerome adds that we who were always in sin shall be saved by
mercy alone. The American English sets the anger and the people's sinfulness
side by side and does not raise the question of order; the explanation
belongs to the Latin the two commentators read.

The next verses go further down. ``All our justices as the rag of a menstruous
woman'', ``we have all fallen as a leaf'', ``There is none that calleth upon
thy name'' and ``thou hast hid thy face from us'' (64:5--6; Douay 64:6--7).
Aquinas reads the hidden face as God's mercy withheld. On the hidden face
Jerome quotes the psalm sung after this reading: \latin{Ostende
faciem tuam, et salvi erimus}, ``shew us thy face, and we shall be saved'' (Ps
80 (79):4), so that we may be servants not of sin but of God (col.~649).

On the last verse, where the people call themselves clay and God the potter,
Jerome turns from their merits to God's clemency. If we consider our merits,
he says, we must despair; if we consider thy clemency, we dare to pray, because thou art
our Father. We are clay, and the pot may not answer back to the potter; but
knowing that we are sons, we dare to say the next verse of the prophet,
\latin{Ne irascaris, Domine, satis}, ``Be not very angry, O Lord''
(cols.~648--650). Aquinas cites Jeremiah~18, where the potter remakes a
spoiled vessel in his hand.

\subsection{Augustine: unless thou convert us}

The psalm's refrain asks in so many words for what Jerome read in Isaiah's
petition: ``Convert us, O God: and shew us thy face, and we shall be saved.'' Augustine's comment on
it begins with one sentence: ``For averse we have been from Thee, and except Thou
convert us, we shall not be converted'' (\work{Enarrationes in Psalmos} 79.3).
Bellarmine says the same in his own way. God does not look on us because we
turn away from him, and ``as we cannot turn to him without his aid, the
prophet asks for both'': that God turn us and that he show his face
(\work{Explanatio in Psalmos} 79, expl.~3).

On the Alleluia verse's psalm Augustine makes the same point at the verse
before it, ``Thou shalt turn us again, and make us alive'' (Douay ``Thou wilt
turn, O God, and bring us to life''): ``not only our being made alive is from
Thee, but our very conversion, that we may be made alive''. On the Alleluia
verse itself he draws the consequence for the one who receives mercy: ``He it
is who cannot indulge in pride, unto whom God showeth His mercy. For by
showing him His salvation He persuadeth him that whatever good man has, he
hath not but from Him who is all our good'' (84.9; §6 in the Oxford
translation). The Collect, too, asks God for the will itself: \latin{Da
\ldots\ hanc tuis fidelibus voluntatem}, the will to go out with just works
to meet Christ.

On the verse that the Communion Antiphon sings, ``the Lord will give
goodness: and our earth shall yield her fruit'', Augustine says where good
fruit comes from. He reads it after the verse ``Truth is sprung out of the earth'', which he
takes as confession: the sinner who confesses his sins makes truth spring
from the earth. Then: \latin{Tua ergo peccata potes habere; fructum bonum non
habebis, nisi ille dederit cui confiteris}, ``Your sins, then, you can have
of your own; good fruit you will not have unless he gives it to whom you
confess'' (84.15). The Lord gives sweetness, ``the sweetness of working
righteousness, so that righteousness shall begin to delight thee, whom before
unrighteousness delighted''; and the preacher may sow, ``and unless God rain
upon it, what profits it that it is sown?'' (§12 in the Oxford translation).

The Entrance Antiphon sings the trust of Psalm 25 (24):2, ``In thee, O my
God, I put my trust; let me not be ashamed.'' Augustine, speaking in the
voice of the Church, explains the trust by what came before it: ``O my God, from
trusting in myself I was brought even to this weakness of the flesh''; and
``now, therefore, `in Thee I trust, I shall not be ashamed''' (24.2).

\subsection{Chrysostom and Aquinas: God is faithful}

Paul thanks God for the Corinthians, says that they lack no gift, and says
that Christ ``will confirm you unto the end without crime''. Chrysostom asks
how Paul can praise a church he is about to call carnal, and answers that
they grew remiss after receiving the gifts, or that Paul speaks of the
greater part as the whole. Whatever the case, he adds, Paul begins with what
only seem to be praises, ``for not even did this praise belong to them, but
to the grace of God''. On v.~8 he notes that ``He shall confirm'' and
``unreprovable'' mark them out ``as still wavering, and liable to reproof'':
they are not yet blameless, and Paul says that Christ will make them so. On
v.~9 Chrysostom gives Paul's reason: ``\,`God is faithful,' i.e. `true.' Now
if `true,' what things He hath promised He will also perform.'' Paul puts
this at the beginning, says Chrysostom, ``lest after the vehemence of the
reproofs they might fall into despair''; and where the call fails,
``unbelief cometh not of Him that calleth, but of those who start away from
Him'' (\work{Homilies on First Corinthians} 2).

Aquinas, on the same verse, states the argument from God's side. God would
not seem faithful \latin{si nos vocaret ad societatem filii et nobis
denegaret, quantum in ipso est, ea, per quae pervenire ad eum possemus}, ``if
he called us to the fellowship of his Son and denied us, as far as lies in
him, the things by which we could reach him''. The fellowship has two times:
\latin{in praesenti per similitudinem gratiae}, ``in the present, by the
likeness of grace'', and \latin{in futuro per participationem gloriae}, ``in
the future, by sharing in glory'' (\work{Super I ad Corinthios} c.~1,
lect.~1).

Bede says something similar of the servants in the Gospel. The master gave
them authority over every work, he says, because by granting the grace of the
Holy Spirit he gave his faithful the power \latin{bonis operibus serviendi},
``to serve in good works'' (\work{In Marci evangelium expositio} IV, on
Mk 13:34). The Prayer over the Offerings offers the gifts \latin{de tuis
\ldots\ collata beneficiis}, from what God has bestowed.

\subsection{Where this interpretation strains}

The first strain is that Isaiah's people ask the question and do not answer
it: the reading does not say why God let them stray. The second is that a
hope resting on God's fidelity can sound like a hope that asks nothing of those who wait. The texts ask a good deal: Isaiah's
people confess, the psalm's people vow not to turn away, the Collect asks for
a will to go out with works, and Chrysostom names the one way the call can
fail: when those who are called ``start away''.

\subsection{The four senses of this interpretation}

\begin{fourSenses}
\item[Literal.] A people that calls the Lord its Father and Redeemer, and
whose sanctuary has been trampled, confesses its sin and appeals to God as
Father and maker (Is 64:7); Paul thanks God for gifts that are God's and
rests the Corinthians' perseverance on God's faithfulness (1~Cor 1:8--9).
The potter and the clay belong to this sense: the prophet's figure for the
people's dependence on their maker.
\item[Allegorical.] The face God is asked to show was veiled with ``a cloud
of flesh'' and was acknowledged when Christ rose and ascended (Augustine on
the refrain, 79.3); those God calls already share the fellowship of his Son
\latin{per similitudinem gratiae}, by the likeness of grace (Aquinas).
\item[Moral.] Humility before the gift: the one to whom God shows mercy
cannot be proud, since every good he has is from God (Augustine 84.9); a man
may have his sins of his own, but good fruit only from the one to whom he
confesses them (84.15); the conversion asked for in the refrain is itself
God's work (79.3). The Prayer after Communion says that by the mysteries God
already teaches his people to love what is heavenly, and the solemn blessing
asks him to make them firm in faith, glad in hope and active in charity.
\item[Anagogical.] Christ will confirm those who wait to the end, without crime
on his day (1~Cor 1:8), and the fellowship had now by the likeness of grace
becomes \latin{participatio gloriae}, a sharing in glory (Aquinas). The
Preface calls what those who keep watch now dare to await a promise
(\latin{promissum}), to be received openly when Christ comes again in glory.
\end{fourSenses}
